\documentclass[11pt,letterpaper]{article}
% Shared layout. Compile each document from this directory.
\usepackage[margin=0.8in,headheight=14pt,headsep=16pt,footskip=22pt]{geometry}
\usepackage[T1]{fontenc}
\usepackage{lmodern,amsmath,amssymb,booktabs,enumitem,titlesec,fancyhdr,microtype,xcolor,hyperref}
\definecolor{inkblue}{RGB}{25,65,100}
\hypersetup{colorlinks=true,linkcolor=inkblue,urlcolor=inkblue,filecolor=inkblue,pdfauthor={AI-assisted course study notes}}
\setlength{\parindent}{0pt}
\setlength{\parskip}{5pt}
\setlength{\emergencystretch}{1.5em}
\setlength{\abovedisplayskip}{7pt}
\setlength{\belowdisplayskip}{7pt}
\titleformat{\section}{\normalsize\bfseries\color{inkblue}}{}{0pt}{}
\titlespacing*{\section}{0pt}{9pt}{3pt}
\setlist[itemize]{leftmargin=15pt,itemsep=2pt,topsep=3pt,parsep=0pt}
\setlist[enumerate]{leftmargin=17pt,itemsep=3pt,topsep=3pt,parsep=0pt}
\pagestyle{fancy}
\fancyhf{}
\fancyhead[L]{\small ECE 750T\enspace /\enspace Responsible AI}
\fancyhead[R]{\small Fall 2026\enspace /\enspace Study notes}
\fancyfoot[R]{\small\thepage}
\renewcommand{\headrulewidth}{0.3pt}
\newcommand{\lesson}[2]{%
  \hypersetup{pdftitle={ECE 750T: #1 -- #2}}%
  {\LARGE\bfseries #1}\par\vspace{3pt}
  {\large #2}\par\vspace{3pt}
  {\footnotesize AI-assisted reconstruction from course slides, not a lecture transcript.
  Discussion reasoning is inferred; classroom responses were not supplied.}\par
  \vspace{3pt}\hrule\vspace{5pt}}
\newcommand{\sources}[1]{\par\vspace{3pt}{\footnotesize\textbf{Slide trail.} #1
  Page numbers refer to PDF pages, not the slides' continuous printed numbering.}\par}
\newcommand{\Prob}{\mathbb{P}}
\newcommand{\E}{\mathbb{E}}
\newcommand{\Yhat}{\widehat{Y}}
\newcommand{\indep}{\mathrel{\perp\!\!\!\perp}}

\begin{document}
\lesson{Lecture 2}{Algorithmic fairness: guarantees, conflicts, and context}

\textbf{The lesson.} Fairness criteria compare different populations and protect
different interests. Choose the comparison explicitly, check its assumptions,
and connect it to the decision's consequences. Equal rates do not establish
individual fairness, appropriate labels, or a fair institutional process.

\section*{Which comparison is being made?}
Let $A$ denote group, $Y$ the recorded binary outcome, $R$ a score, and $\Yhat$
the binary decision. Discrimination can affect allocation, quality of service,
or representation. Excluding $A$ from inputs leaves proxies and biased labels:
the slides' school-based rule awards 80\% of one group and 20\% of another
despite never reading the group label.

\begin{center}
\small
\begin{tabular}{@{}p{0.25\linewidth}p{0.70\linewidth}@{}}
\toprule
Criterion & What is equal across groups? \\
\midrule
Demographic parity & $\Prob(\Yhat=1\mid A)$; independence: $\Yhat\indep A$. \\
Equal opportunity & TPR: $\Prob(\Yhat=1\mid Y=1,A)$. \\
Equalized odds & TPR and FPR; separation: $\Yhat\indep A\mid Y$. \\
Predictive parity & PPV: $\Prob(Y=1\mid\Yhat=1,A)$. \\
Sufficiency & $Y\indep A\mid R$; binary decisions require equal PPV \emph{and} NPV. \\
Group calibration & $\Prob(Y=1\mid R=r,A=a)=r$ for every supported score and group. \\
\bottomrule
\end{tabular}
\end{center}

Calibration overall can conceal group miscalibration: the slides' score 0.7 has
outcome rates 0.8 and 0.6 in two groups, averaging to 0.7. Group calibration
implies score sufficiency where groups share support. Thresholding a calibrated
score need not preserve predictive parity: groups can have different mixtures
of scores above the cutoff. Condition only on events with positive probability.

\section*{Two conflicts worth being able to derive}
Write $\pi_a=\Prob(Y=1\mid A=a)$ and $q_a=\Prob(\Yhat=1\mid A=a)$. Then
\[
q_a=\mathrm{TPR}_a\pi_a+\mathrm{FPR}_a(1-\pi_a).
\]
Under equalized odds, with common TPR $u$ and FPR $v$,
$q_a-q_b=(u-v)(\pi_a-\pi_b)$. Unequal base rates and $u\ne v$ prevent
demographic parity. The assumptions matter: $u=v$ is an exception, and the
binary-outcome argument should not be extended blindly to multiclass targets.

Bayes' rule gives the second conflict:
\[
\mathrm{PPV}_a=
\frac{\mathrm{TPR}_a\pi_a}
{\mathrm{TPR}_a\pi_a+\mathrm{FPR}_a(1-\pi_a)}.
\]
With $0<\pi_a,\pi_b<1$, unequal base rates and common \emph{positive} TPR and
FPR imply unequal PPV. Perfect prediction is an important exception. The COMPAS
discussion uses this tension to explain why predictive parity and equal error
rates can disagree, without deciding which criterion a court should prioritize.

Errors acquire consequences through policy. If a high-risk label changes a penalty
from $t_{\min}$ to $t_{\max}$, the expected penalty for recorded $Y=0$ in group
$a$ is $t_{\min}+(t_{\max}-t_{\min})\mathrm{FPR}_a$. Using the same score
to offer support changes the meaning of that false positive.

\newpage
\section*{Individual fairness moves the hard choice into a metric}
Dwork et al. model a randomized decision as $M(x)\in\Delta(O)$ and require
\[
D_{\mathrm{TV}}(M(x),M(x'))\le d(x,x'),\qquad
D_{\mathrm{TV}}(P,Q)=\tfrac12\sum_o|P(o)-Q(o)|.
\]
For binary decisions with award probabilities $p_x,p_{x'}$, total variation is
$|p_x-p_{x'}|$. With $d(x,x')=0.1$, probabilities $0.8,0.3$ violate the
constraint; $0.55,0.45$ satisfy it. The guarantee concerns distributions over
treatment, not equality of every realized draw.

Equal group award rates can still treat similar cross-group pairs differently.
Conversely, individual fairness permits a group gap if groups differ under the
chosen task metric. If $S,T$ are distributions over individuals and $\mu_S,\mu_T$
their induced outcome distributions, then
\[
D_{\mathrm{TV}}(\mu_S,\mu_T)\le d_{\mathrm{EM}}(S,T),\qquad
d_{\mathrm{EM}}(S,T)=\min_{\gamma\in\Pi(S,T)}\E_\gamma[d(x,x')].
\]
Earthmover distance is the cheapest average task distance for transporting one
group distribution to the other. Close groups force close outcome distributions.
For the slides' groups at $\{.1,.3\}$ and $\{.6,.8\}$ with equal masses,
$d_{\mathrm{EM}}=.5$: the fair rule $p_x=x$ attains gap .5, whereas $p_x=.5$
has gap zero. A bound permits a disparity; it does not require it.
The metric's relevance, scale, and cross-group comparisons need justification.

\section*{Where interventions act, and what A1 adds}
\textbf{Preprocessing} changes data, weights, labels, samples, or representations.
Adversarial representation learning trades task information against the ability
to infer group membership; synthetic-data methods also impose assumptions.
\textbf{In-processing} adds constraints or penalties to training.
\textbf{Post-processing} changes decisions, for example via thresholds or rejection
options. Feasibility, access to group information, utility, and transfer to new
settings depend on the method; the slides' comparison is illustrative.

A1's smooth penalty can use
$(\overline{R}_{A=0}-\overline{R}_{A=1})^2$ alongside prediction loss.
It supplies gradients, unlike hard thresholded decisions, but score-mean parity
does not guarantee binary-decision parity. Evaluate the hard DP gap separately.
Reweighted accuracy averages the two group accuracies equally; ordinary accuracy
weights by group frequency. Compute evaluation rates on the full evaluation set.

\section*{The five abstraction traps explain what a metric omits}
\textbf{Framing:} a narrow model boundary omits how outputs become decisions.
\textbf{Portability:} social context changes what a method means, beyond a change
in data distribution. \textbf{Formalism:} a mathematical condition captures only
part of fairness. \textbf{Ripple effect:} deployment changes behaviour and
institutions. \textbf{Solutionism:} beginning with a model excludes policy changes
or not deploying. A criterion alone cannot establish legal compliance.

\textbf{Midterm emphasis (inference).} Highest priority: compute and interpret
group rates, distinguish conditioning directions, and derive conflicts with their
exceptions. Also practise the individual-fairness bound, intervention choices, and
applying an abstraction trap to a concrete case. The reference sheet supplies the
main group criteria, but not the individual-fairness or earthmover formulas.

\sources{\href{../L02.pdf}{L02.pdf}, pp.~6--18 (group criteria), 19--30
(interventions/A1), 31--45 (Dwork et al.), 46--55 (Chouldechova),
56--65 (Selbst et al.); annotated diagrams consulted. A1: pp.~1--3.}
\end{document}
